\usepackage[colorlinks,linkcolor=black,urlcolor=black]{hyperref}
\usepackage{amsmath,amssymb,amsthm}
\usepackage[utf8]{inputenc}
\usepackage{xcolor}
\usepackage{xparse}
\usepackage{style}
\usepackage[capitalise]{cleveref}
\usepackage{mathtools}
\usepackage[shortlabels]{enumitem}
\usepackage[style=apa, backend=biber]{biblatex}

\usepackage[record,style=alttreegroup,nomain,symbols]{glossaries-extra}

\usepackage{dutchcal} 

\usepackage{pdfpages}

\usepackage[T1]{fontenc}
%%%%%%%%%%%%%%%%
% Seitenlayout %
%%%%%%%%%%%%%%%%

% DIV# gibt den Divisor für die Layoutberechnung an.
% Vergrößern des Divisors vergrößert den Textbereich.
% BCOR#cm gibt die Breite des Bundstegs an.
\usepackage[DIV14,BCOR2cm]{typearea}

% Abstand obere Blattkante zur Kopfzeile ist 2.54cm - 15mm
%\setlength{\topmargin}{-15mm}

% Keine Einrueckung nach einem Absatz.

%\parindent 0pt

% Abstand zwischen zwei Abs\"atzen.

%\parskip 12pt

% Zeilenabstand.

\linespread{1.25}



\usepackage{makeidx}
\makeindex
